\section{论文的主要内容和技术路线}

\subsection{主要研究内容}

本研究围绕 GPU 驱动的渲染管线进行探讨，重点关注流式加载与 LOD 机制在工业级大规模高精度 CAD 模型渲染中的应用。

\subsubsection{GPU驱动的渲染管线（GPU-Driven Pipeline）}

传统的 CPU 驱动渲染管线依赖 CPU 完成剔除计算、 LOD 选择等任务，随着工业场景模型复杂度的增加， CPU 端的计算开销逐渐成为瓶颈。 GPU 驱动的渲染管线通过将这些计算任务转移至 GPU 端执行，有效提升了渲染效率，并减少了 CPU 负载。

本研究主要探讨 GPU 驱动渲染管线的关键技术，包括：

（1）基于簇的渲染（Cluster-Based Rendering）：将几何数据划分为更小的簇（Cluster），并利用任务着色器（Task Shader）+ 网格着色器（Mesh Shader）的组合取代传统的计算着色器（Compute Shader）+ 顶点着色器（Vertex Shader） + 几何着色器（Geometry Shader）方案，实现高效的簇级剔除、LOD选择和渲染\cite{Czubala2024}。

（2）GPU 端剔除计算：利用计算着色器或任务着色器在 GPU 上进行视锥剔除（Frustum Culling）、遮挡剔除（Occlusion Culling）和圆锥剔除（Cone Culling），避免渲染不可见或不必要的几何数据；

（3）LOD 选择：结合屏幕空间误差及 GPU 并行计算能力，在 GPU 端动态选择适合的 LOD 级别，降低渲染开销。

\subsubsection{流式加载（Streaming）}

对于大规模工业模型而言，一次性加载整个场景会导致显存溢出或不必要的资源占用，而流式加载机制可以根据模型的可见性和重要性按需加载资源，提高显存利用率。本研究将重点探讨以下流式加载关键技术：

（1）页表（Page Table）：建立多级页表结构，并基于访问频率、视角变化、 GPU 负载情况进行智能换入换出策略设计，减少不必要的数据加载，提高资源利用率；

（2）异步加载（Async Update）：进行后台数据加载，减少数据加载对主渲染流程的影响，保证渲染的实时性；

（3）编码和解码（Encode and Decode）：采用高效的几何压缩算法对模型数据进行存储，以减少传输和存储开销。在 GPU 端，通过并行解码技术实现高效的实时解压缩，确保数据在加载后能迅速投入渲染，同时减少对显存带宽的占用，提高数据加载效率\cite{Mlakar2024}；

（4）优先级调度（Priority Scheduling）：基于模型的重要性（如视距远近、屏幕占比、遮挡关系等）对加载任务进行优先级排序，使关键模型数据优先加载，以提升画面质量与交互流畅度。

\subsubsection{LOD（Level of Detail）机制}

LOD 机制是优化大规模场景渲染的关键技术之一。通过降低远距离或非关键区域的几何复杂度，可以有效减少 GPU 计算量，从而提升渲染性能\cite{Deng2017}。本研究将探讨以下关键内容：

（1）LOD 生成（LOD Generation）：分析适用于工业模型的 LOD 生成方法，包括基于网格简化（Mesh Simplification）和基于簇合并的 LOD 生成\cite{Jensen2023}。同时，为了避免 LOD 切换时的裂缝和视觉伪影，可采用拓扑约束的网格简化算法，确保不同 LOD 级别间的边界对齐；

（2）运行时 LOD 选择（Runtime LOD Selection）：研究基于屏幕像素误差（Screen Space Error, SSE）的动态 LOD 选择策略，确保在视觉质量与性能之间取得平衡\cite{WangQian2016}。

\subsection{技术路线}

本研究采用 Vulkan 作为图形 API ，并基于开源自研引擎进行开发。

Vulkan 作为一款低级别、高性能的图形 API ，支持GPU端任务调度、细粒度内存管理和跨平台部署，能够有效提升渲染效率\cite{overvoorde2020vulkan}。其显式控制的资源管理模式减少了 CPU 开销，使 GPU 驱动渲染管线成为可能。此外， Vulkan 提供稀疏资源（Sparse Resources）管理，可用于实现高效的流式加载，优化超大规模数据的显存利用率。

自研引擎的开放架构便于扩展流式加载、基于簇的渲染、 LOD 机制等关键技术，从而优化超大规模数据的渲染效率。

接下来是整体技术路线：

（1）预处理阶段：

在模型导入内存后，将进行 LOD 生成，并为每个簇构建相应的顶点数据与索引数据。参考 Nanite 的处理方式，簇的划分采用 MeshOptimizer 库，而簇组（Meshlet/Cluster Group）的划分则基于 METIS 库，以确保在 LOD 过渡时实现无缝切换，避免渲染过程中可能出现的视觉伪影。

为了提高簇组简化过程的稳定性，可借助 KD-Tree 进行顶点合并，从而优化网格结构。此外，为了降低预处理阶段的计算开销并提升处理效率，在数据量较大的情况下，可利用 oneTBB 库进行多线程优化，以充分发挥多核处理能力，加速 LOD 构建与簇组划分的整体流程。

（2）实例级剔除（Instance Culling）和 LOD 选择：

在 GPU 端通过计算着色器进行实例级剔除（包括视锥剔除等），利用提前计算的每个实例的包围盒信息及所有簇的屏幕像素误差，决定哪些实例的簇需要加载至显存。此过程能够有效减少冗余数据的传输和存储，提高渲染效率。

（3）流式加载：

根据计算着色器的剔除与 LOD 选择结果，确定需要加载至显存的页面（Page）。所有页面使用页表进行管理，并在显存资源受限时执行数据换入换出（Swap In/Out）以维持高效的显存利用率。此部分将基于 Vulkan 的稀疏资源机制，以实现更灵活的显存管理，从而支持大规模数据的高效调度。

（4）最终绘制：

在网格着色器中执行簇级别的视锥剔除和圆锥剔除，随后依据 LOD 规则动态调整渲染级别。只有通过剔除和 LOD 筛选的簇才会进入最终的渲染管线，以最大程度优化渲染性能，实现高效的实时绘制。


\subsection{可行性分析}

本研究结合 Vulkan 的高性能渲染特性与自研引擎的可扩展架构，在技术上具备较高的可行性。而 Nanite 作为 GPU 驱动渲染管线的成功案例，证明了基于 GPU 的 LOD 管理、剔除计算及基于簇的渲染方案的可行性。近年来， Vulkan 生态系统不断完善，已有众多成熟的稀疏绑定（Sparse Binding）、网格着色器、虚拟纹理（Virtual Texture）等相关实践，可为本研究提供技术借鉴。